%\documentclass[12pt]{amsart}
\usepackage{amscd}
\usepackage{amsmath,amsthm}     % Handy math stuff, theorem environments.
\usepackage{amssymb}
\usepackage{euscript}           % Nice script font.
\usepackage[all,2cell,ps]{xy} \UseAllTwocells \SilentMatrices
%\usepackage{makeidx}
\CompileMatrices		% speeds up XY-pic diagram compilation
\UseComputerModernTips		% use Knuth's computer modern fonts

% Comment out before distribution - 
%\usepackage{showkeys}  % Enables the labels in text.
%\usepackage{hyperref}  % Hyperlinks

% Comment in for smaller page margin
%\addtolength{\oddsidemargin}{-1cm}
%\addtolength{\evensidemargin}{-1cm}
%\addtolength{\headheight}{0pt}
%\addtolength{\headsep}{0cm}
%\addtolength{\textheight}{0cm}
%\addtolength{\textwidth}{2cm}
%\addtolength{\footskip}{0cm}

%\parskip1ex 


%%% This part of the file (after the \documentclass command,
%%% but before the \begin{document}) is called the ``preamble''.
%%% This is a good place to  put our macro definitions.

\newcommand{\tensor}{\otimes}
\newcommand{\homotopic}{\simeq}
\newcommand{\homeq}{\cong}
\newcommand{\iso}{\approx}

\DeclareMathOperator{\ho}{Ho}

\newcommand{\R}{\mathbb{R}}
\newcommand{\C}{\mathbb{C}}
\newcommand{\Z}{\mathbb{Z}}

\newcommand{\M}{\mathcal{M}}
\newcommand{\W}{\mathcal{W}}

\newcommand{\itilde}{\tilde{\imath}}
\newcommand{\jtilde}{\tilde{\jmath}}
\newcommand{\ihat}{\hat{\imath}}
\newcommand{\jhat}{\hat{\jmath}}
\newcommand{\comment}[1]{}

\newcommand{\cat}[1]{\EuScript {#1}}
\newcommand{\axiom}[1]{\textbf {#1}}

% arrows
\newcommand{\lra}{\longrightarrow}
\newcommand{\lla}{\longleftarrow}
\newcommand{\ra}{\rightarrow}
\newcommand{\la}{\leftarrow}
\newcommand{\llra}{\longleftrightarrow}
\newcommand{\dbra}{\rightrightarrows}           % double arrow for equalizer diagrams
\newcommand{\Ra}{\Rightarrow}
\newcommand{\La}{\Leftarrow}

% dots
%\newcommand{\tdot}{^{^\centerdot}} 	% upper dot
\newcommand{\tdot}{^{^{_{_\bullet}}}} 	% upper dot
\newcommand{\tdotdot}{^{^{_{_{\bullet\bullet}}}}}
%\newcommand{\sdot}{_\centerdot}			% lower dot
\newcommand{\sdot}{_{_{^{^\bullet}}}} 	% upper dot
\newcommand{\sdotdot}{_{_{^{^{\bullet\bullet}}}}} 	% upper dot

%math symbols
\newcommand{\comp}{\circ}

%special faces & letters
\newcommand{\eps}{\epsilon}
\newcommand{\epso}{\overline \epsilon}
\newcommand{\nuo}{\overline \nu}
\newcommand{\Deltao}{\overline \Delta} 
\newcommand{\Mo}{\overline M} 

%completion
\newcommand{\compl}[1]{{#1}\,{\hat{}}_{\!\!\!R}}

%%%%%%%%%%%%%%%%%%%%%%%%%%%%%%%%%%%%%%%%%%%%%%%%%%%%%%%%%%%%%%%%%%%%%%%%%%%%%%
\newdir{ >}{{}*!/-5pt/\dir{>}}                  % Make better tailed arrows.
\newdir{ |}{{}*!/-5pt/\dir{|}}                  % Make better mapsto arrows.

% Some theorem-like environments, all numbered together starting at 1
% in each section.

\newcommand{\mdfn}[1]{\dfn{\mathversion{bold}#1}} % Even make math symbols bold

% Topologists smash /\ and wedge \/.
\newcommand{\Smash}             {\wedge}
\newcommand{\bigSmash}          {\bigwedge}
\newcommand{\hfsm}              {\rtimes}
\newcommand{\eqhsma}[1]{\underset{\Sigma_{#1}}  %`equivariant half-smash'---
                                 {\rtimes}}     %\eqhsma{n} produces the symbol
                                                %`half-smash over sigma_n'
\newcommand{\eqtens}[1]{\underset{\Sigma_{#1}}  %`equivariant tensor' 
                                 {\otimes}}
% For fibred products. Modify this to something fancier if desired.
% "arg1 multiplied over arg2 with arg3"
\newcommand{\fibpr}[3]{{#1}\times_{#2}{#3}}

% produces a times symbol with a subscript directly beneath it, rather
% than beneath and to the right.
\newcommand{\timessub}[1]{\underset{#1}{\times}}

% same, but with otimes
\newcommand{\otimessub}[1]{\underset{#1}{\otimes}}

\newcommand{\und}{\underline}
\newcommand{\norm}[1]{\mid #1 \mid}
\newcommand{\Sum} {\sqcup}
\newcommand{\Wedge} {\vee}
\DeclareRobustCommand{\bigWedge}{\bigvee}
\DeclareMathOperator*{\invlim}{\varprojlim}
\DeclareMathOperator*{\dirlim}{\varinjlim}
\DeclareMathOperator*{\colim}{colim\,}
\DeclareMathOperator*{\wcolim}{wcolim\,}
\DeclareMathOperator*{\wlim}{wlim\,}
\DeclareMathOperator*{\initial}{{\bf 0}}
\DeclareMathOperator*{\final}{{\bf 1}}
\DeclareMathOperator*{\terminal}{\final}
\DeclareMathOperator*{\initialcat}{\emptyset}
\DeclareMathOperator*{\finalcat}{{\it e}}
\DeclareMathOperator*{\terminalcat}{\finalcat}

\newcommand{\corollaryref}[1] {Cor. \ref{#1}}
\newcommand{\definitionref}[1] {Def. \ref{#1}}
\newcommand{\lemmaref}[1] {Lemma \ref{#1}}
\newcommand{\propositionref}[1] {Prop. \ref{#1}}
\newcommand{\theoremref}[1] {Thm. \ref{#1}}
\newcommand{\conjectureref}[1] {Conj. \ref{#1}}
\newcommand{\remarkref}[1] {Remark \ref{#1}}
\newcommand{\chapterref}[1] {Chap. \ref{#1}}
\newcommand{\sectionref}[1] {Section \ref{#1}}
\newcommand{\equationref}[1] {(\ref{#1})}


% The following causes equations to be numbered within sections
%\newcounter{theoremcounter}
%\numberwithin{theoremcounter}{subsection}
%\numberwithin{equation}{section}

%\theoremstyle{plain} %% This is the default, anyway
%\newtheorem{thm}[theoremcounter]{Theorem}
%\newtheorem{cor}[theoremcounter]{Corollary}
%\newtheorem{lem}[theoremcounter]{Lemma}
%\newtheorem{prop}[theoremcounter]{Proposition}
%\newtheorem{conjecture}[theoremcounter]{Conjecture}
%\theoremstyle{definition}
%\newtheorem{defn}[theoremcounter]{Definition}
%\theoremstyle{remark}
%\newtheorem{rem}[theoremcounter]{Remark}
%\newtheorem{ex}[theoremcounter]{Example}
%\newtheorem{notation}[theoremcounter]{Notation}
%\newtheorem{terminology}[theoremcounter]{Terminology}
